This research addresses a specific methodological gap in digital media creation by developing systematic approaches for translating cultural epistemologies into computational procedures within interactive systems. The primary contribution is Speculocultural Technopoiesis as a documented methodology that demonstrates how cultural practices can be encoded into algorithmic frameworks while maintaining community control over the cultural translation process.

\textbf{Methodological Contribution:} The research provides concrete protocols for cultural-technical translation in digital media creation, offering documented procedures for identifying cultural assumptions in existing computational systems and systematically implementing community-controlled modifications. Through cyclical processes of Unsettling, Remixing, Embodying, and Transmitting, ST offers practical techniques that other cultural communities can adapt for their own digital media projects while maintaining cultural authenticity and community accountability.

\textbf{Practical Demonstration:} Three interactive works created through ST methodology document specific approaches to encoding Black epistemologies into computational procedures, providing proof-of-concept examples of how cultural practices translate into interaction mechanics and procedural frameworks. These works serve as concrete demonstrations of the methodology's application rather than theoretical abstractions, offering documented techniques that practitioners can study and adapt.

\textbf{Community-Centered Process:} By positioning community accountability as integral to technical implementation rather than external oversight, the research contributes practical models for community-controlled digital media creation that avoid both extractive research practices and technical fixes that ignore cultural sovereignty. The methodology demonstrates how communities can systematically appropriate and modify existing digital creation tools according to their own epistemological frameworks.